% ============================================================
% §3  Three distance regimes
% ============================================================
\section{Three interaction regimes}\label{sec:rezimy}

The central prediction of TGP is the existence of \textbf{three
distinct regimes} of interaction, arising from the interference of the
$\Phi$~fields generated by different sources.  These are not three
separate forces --- they are \textbf{one dynamics} of the field~$\Phi$
that yields a different sign of the net force at different distance
scales.

% ----------------------------------------------------------
\subsection{Mechanism: gradient of $\Phi$ as force}\label{ssec:mechanizm}

A test object moving in the field~$\Phi$ experiences a force
proportional to the gradient:
\begin{equation}\label{eq:force}
    \mathbf{F} \;\propto\; -\nabla V_{\mathrm{eff}}(\Phi),
\end{equation}
where $V_{\mathrm{eff}}(\Phi)$ is the effective potential felt by a
particle in the spatial field.

Key observation: $V_{\mathrm{eff}}$ is \textbf{not a monotonic}
function of the distance from the source.  This follows from the
nonlinear coupling $\mathcal{N}[\Phi]$
(Eq.~\ref{eq:Phi-nonlinear}).  The effective potential between two
sources has the form

\begin{equation}\label{eq:Veff-profile}
    V_{\mathrm{eff}}(d) \;=\;
    \underbrace{-\frac{(q\PhiZero)^2 M_1 M_2}{4\pi\,d}}_{%
      E_{\mathrm{lin}}:\;\text{attraction (large }d\text{)}}
    \;+\;
    \underbrace{\frac{2\beta\,(q\PhiZero)^2 M_1 M_2}
      {(4\pi)^2\,\PhiZero\,d}\ln\frac{d}{r_0}}_{%
      E_\beta:\;\text{repulsion (intermediate }d\text{)}}
    \;-\;
    \underbrace{E_\gamma(d)}_{%
      \sim d^{-3}:\;\text{well (small }d\text{)}}
\end{equation}
where $E_\gamma(d) > 0$ is a quartic term that dominates at
small~$d$.  The crucial feature is the \textbf{structure}: three
sign changes of the gradient, resulting from the different growth
rates of the linear ($\sim 1/d$), cubic (repulsive), and quartic
(well) terms as the sources approach each other.

\begin{remark}[Explicit link to the field equation]\label{rem:Veff-field-eq}
The terms in the potential~\eqref{eq:Veff-profile} follow directly
from the field equation~\eqref{eq:full-field-alt}:
\begin{itemize}[nosep]
  \item $E_{\mathrm{lin}} \sim -1/d$: from the linear part
    $\nabla^2\Phi/\PhiZero$ after double volume integration between
    two Gaussian spheres;
  \item $E_\beta \sim +1/d^2$ (repulsion): from the cubic term
    $\beta\Phi^2/\PhiZero^2$ in $\mathcal{N}[\Phi]$; dominant when
    gradients and the value of $\Phi$ are moderately large;
  \item $E_\gamma \sim -1/d^3$ (well): from the quartic term
    $-\gamma\Phi^3/\PhiZero^3$; dominant at very large $\Phi$
    (extremely close sources, Regime~III).
\end{itemize}
The exact decomposition
(Eqs.~\ref{eq:Eint-decomp}--\ref{eq:Egamma})
is given in \S\ref{ssec:dwa-zrodla}.
\end{remark}

% ----------------------------------------------------------
\subsection{Regime I: large scales --- gravity}\label{ssec:grawitacja}

\begin{proposition}[Gravity as an interference effect]\label{prop:grav}
At large scales ($r \gg r_{\mathrm{rep}}$, where $r_{\mathrm{rep}}$
is the repulsion scale) the interference of the $\Phi$~fields
generated by two material objects causes the gradient of~$\Phi$ to
point \textbf{toward the other object}.  The objects are pushed
together.
\end{proposition}

Mechanism: each object produces a ``hill'' of $\Phi$ around itself.
Two distant objects have between them a region where their fields
overlap \emph{constructively} --- $\Phi$ between them is higher than
on the outside.  A test object ``rolls'' down the $\Phi$~gradient
toward the other body.

This is \textbf{gravity} in TGP: not curvature of a pre-existing
geometry, but a force arising from the constructive interference of
generated space.

In the weak-field, large-distance limit one recovers Newton's law:
\begin{equation}\label{eq:newton-limit}
    F_{\mathrm{grav}} \simeq
    -\frac{G_{\mathrm{eff}}\,M_1\,M_2}{r^2}
    \qquad (r \gg r_{\mathrm{rep}}).
\end{equation}
The constant $G_{\mathrm{eff}}$ is here a \textbf{derived quantity},
expressed in terms of the generation constant~$q$ and the parameters
of the profile $\phi(r)$, not introduced by hand.

% ----------------------------------------------------------
\subsection{Regime II: intermediate scales --- repulsion}%
\label{ssec:odpychanie}

\begin{proposition}[Spatial repulsion]\label{prop:rep}
At intermediate scales
($r_{\mathrm{well}} < r < r_{\mathrm{rep}}$) the spatial potential
near particles is so high that the gradient of~$\Phi$ points
\textbf{away from the source}.  Objects are repelled.
\end{proposition}

Mechanism: in the immediate neighbourhood of a particle the
field~$\Phi$ attains very high values (the particle generates dense
space around itself).  When a second object approaches, it enters a
region of \emph{steep rise} in~$\Phi$ --- the gradient points
outward, meaning a repulsive force.

This repulsion may account for:
\begin{itemize}[nosep]
  \item stability of matter (particles do not collapse onto each
        other),
  \item quantum effects (the uncertainty principle as a consequence
        of spatial repulsion at small scales),
  \item degeneracy pressure (a many-body effect in confined volume).
\end{itemize}

\begin{remark}[Connection with the uncertainty principle]
In standard quantum mechanics the Heisenberg uncertainty principle
($\Delta x \cdot \Delta p \geq \hbar/2$) is an axiom.  In TGP it
may \emph{follow} from spatial repulsion: attempting to localise a
particle more tightly means entering a region of high~$\Phi$ and
strong gradient, which generates a ``kick-back'' --- a physical
source of uncertainty.  The following proposition formalises this
mechanism.
\end{remark}

\begin{proposition}[Uncertainty principle from TGP spatial repulsion]%
\label{prop:uncertainty-tgp}
Let a particle of mass~$m$ be localised within a region of
size~$\Delta x$ in the field $\Phi(r)$ generated by itself
(self-coupling, Axiom~\ref{ax:zrodlo}).  In Regime~II (repulsion,
$r_{\mathrm{well}} < \Delta x < r_{\mathrm{rep}}$) the field energy
in the core volume $V_{\rm core} \sim (4\pi/3)\Delta x^3$ is
\begin{align}\label{eq:E-field-core}
  E_{\rm field}(\Delta x)
  &= \int_0^{\Delta x} 4\pi r^2\,
     \frac{\beta\,\Phi_{\rm loc}^2(r)}{\PhiZero^2}\,dr
  \;\approx\;
  \frac{\beta\,(qm)^2}{4\pi\,\PhiZero^2}\,\Delta x,
\end{align}
where $\Phi_{\rm loc}(r) = qm/(4\pi r)$ is the field in the linear
approximation.
\textbf{Key point}: the field energy grows linearly with $\Delta x$
(it accumulates in an ever-larger core).  The kinetic energy from
momentum uncertainty $\Delta p \sim \hbar(\Phi)/\Delta x$ is
\begin{equation}\label{eq:E-kin-uncertainty}
  E_{\mathrm{kin}} \sim \frac{(\Delta p)^2}{2m}
  = \frac{\hbar(\Phi)^2}{2m\,(\Delta x)^2}.
\end{equation}
Minimising
$E_{\mathrm{total}} = E_{\mathrm{kin}} + E_{\rm field}$ with
respect to~$\Delta x$:
\begin{equation}\label{eq:Emin-condition}
  \frac{dE_{\rm total}}{d(\Delta x)} = 0
  \quad\Longrightarrow\quad
  -\frac{\hbar(\Phi)^2}{m\,(\Delta x^*)^3}
  + \frac{\beta\,(qm)^2}{4\pi\,\PhiZero^2} = 0,
\end{equation}
yielding the minimal localisation scale
\begin{equation}\label{eq:Dx-min}
  \Delta x^* = \left(
    \frac{4\pi\,\PhiZero^2\,\hbar^2(\Phi)}
         {m\,\beta\,(qm)^2}
  \right)^{\!1/3}.
\end{equation}
The corresponding momentum is
$\Delta p^* = \hbar(\Phi)/\Delta x^*$, so the product equals
\begin{equation}\label{eq:uncertainty-tgp-full}
  \Delta x^*\cdot\Delta p^* = \hbar(\Phi).
\end{equation}
Hence the uncertainty bound:
\begin{equation}\label{eq:HUP-TGP}
  \boxed{\Delta x \cdot \Delta p
  \;\geq\; \hbar(\Phi)
  = \hbar_0\sqrt{\PhiZero/\Phi}.}
\end{equation}
For $\Phi = \PhiZero$ (reference conditions):
$\hbar(\PhiZero) = \hbar_0$, recovering the standard Heisenberg
uncertainty principle \textbf{as a dynamical result}, not an axiom.

\medskip\noindent\emph{Status and scope of validity.}\;
The derivation~\eqref{eq:E-field-core}--\eqref{eq:Dx-min} is a
\textbf{schematic order-of-magnitude proof} for the linear-field
model ($\Phi_{\rm loc} \sim qm/(4\pi r)$) in the weak-field regime.
Full precision requires: (a)~spinorial connection in TGP
(Section~\ref{sec:formalizm}), (b)~quantisation of~$\Phi$ with the
substrate cutoff (Appendix~\ref{app:kwantyzacja}).
\end{proposition}

\begin{corollary}[Variability of the uncertainty bound with $\Phi$]%
\label{cor:uncertainty-variable}
In a strong field ($\Phi \gg \PhiZero$, near a compact object):
$\hbar(\Phi) = \hbar_0\sqrt{\PhiZero/\Phi} \ll \hbar_0$.
The effective uncertainty bound is \textbf{suppressed}:
\[
  \Delta x \cdot \Delta p \;\geq\; \hbar(\Phi)/2
  \;\ll\; \hbar_0/2.
\]
Matter in a strong field is \emph{closer to classical}: quantumness
fades naturally without additional axioms.
\end{corollary}

% ----------------------------------------------------------
\begin{proposition}[Substrate propagation speed $v_W(\chi)$]%
\label{prop:vW-substrate}
Let $\chi \equiv \Phi/\Phi_0$ be the normalised spatial scalar.
The covariant field equation of TGP,
\begin{equation}\label{eq:field-eq-covariant}
  \Box_g\Phi \;=\; \mathcal{N}[\Phi],
\end{equation}
linearised around a homogeneous background $\chi = \chi_0$, yields
the substrate perturbation speed
\begin{equation}\label{eq:vW-formula}
  v_W(\chi)
  \;=\; \frac{c_0}{\sqrt{\chi}},
\end{equation}
where $c_0$ is the speed of light in the reference state
($\chi = 1$).  In particular: $v_W(1) = c_0$,
$v_W(\chi) \to \infty$ as $\chi \to 0$ (vacuum region),
$v_W \to 0$ as $\chi \to \infty$
(interior of $\mathcal{S}_\infty$).

\medskip\noindent
\emph{Tensor GW speed}:
$v_{\mathrm{GW}}^{\mathrm{tensor}} = c_0$ (the $\sigma_{ab}$
sector, independent of~$\chi$) --- consistent with the observation
GW170817.
\end{proposition}

% ----------------------------------------------------------
\subsection{Regime III: extremely small scales --- confining well}%
\label{ssec:studnia}

\begin{proposition}[Spatial well]\label{prop:well}
At extremely small scales ($r < r_{\mathrm{well}}$) the nonlinear
self-coupling $\mathcal{N}[\Phi]$ creates a \textbf{potential well}
--- a region in which two sources are ``glued'' together in a common
trough of the field~$\Phi$.  The gradient of~$\Phi$ once again
points inward.
\end{proposition}

Mechanism: when two sources are extremely close, their $\Phi$~fields
overlap so strongly that the nonlinear terms $\mathcal{N}[\Phi]$
dominate and create a deep well --- a region from which escape
requires supplying significant energy.

This regime may account for:
\begin{itemize}[nosep]
  \item the strong nuclear interaction (quark confinement): quarks
        are ``glued'' together in a shared spatial well;
  \item formation of hadronic bound states;
  \item asymptotic freedom: inside the well the gradients are small
        (flat bottom), so particles behave as though free.
\end{itemize}

% ----------------------------------------------------------
\subsection{Matter as defects of the field $\Phi$}\label{ssec:defekty}

In TGP a material particle is a \textbf{localised defect} of the
field~$\Phi$ --- a radial kink.

\begin{definition}[Radial kink]\label{def:kink}
A \emph{radial kink} of the field~$\Phi$ is a spherically symmetric,
localised solution of
\begin{equation}\label{eq:kink}
    \nabla^2\Phi - V'_{\mathrm{eff}}(\Phi) = 0,
\end{equation}
satisfying the boundary conditions
\begin{equation}\label{eq:kink-bc}
    \Phi(0) = \Phi_{\!0} \neq \PhiZero, \qquad
    \Phi(r\to\infty) = \PhiZero.
\end{equation}
\end{definition}

\begin{proposition}[Defect mass]\label{prop:masa-defektu}
The mass associated with a radial kink of~$\Phi$ is
\begin{equation}\label{eq:masa-defektu}
    m_{\mathrm{def}} = 4\pi \int_0^{\infty}
    \left[\,\frac{1}{2}\!\left(\frac{d\Phi}{dr}\right)^{\!2}
    + V_{\mathrm{eff}}(\Phi)
    - V_{\mathrm{eff}}(\PhiZero)\right] r^2\,dr.
\end{equation}
The integral is finite owing to the boundary condition
$\Phi(r\to\infty) = \PhiZero$ and regular behaviour of $\Phi$ at the
origin.
\end{proposition}

% ----------------------------------------------------------
\subsection{Summary: effective-potential diagram}\label{ssec:diagram}

The effective potential $V_{\mathrm{eff}}(r)$ between two sources
exhibits the following structure:

\begin{center}
\begin{tabular}{ccc}
\toprule
\textbf{Scale} & $\nabla V_{\mathrm{eff}}$ & \textbf{Effect} \\
\midrule
$r < r_{\mathrm{well}}$ & $\to$ inward & Well (confinement) \\
$r_{\mathrm{well}} < r < r_{\mathrm{rep}}$
  & $\leftarrow$ outward & Repulsion \\
$r > r_{\mathrm{rep}}$ & $\to$ inward & Gravity (attraction) \\
\bottomrule
\end{tabular}
\end{center}

The shape of $V_{\mathrm{eff}}$ resembles a Lennard-Jones potential
with an additional well at small scales --- but its origin is
fundamentally different: it arises not from empirical fitting, but
from the nonlinear interference of the field~$\Phi$.
